% ---
% Reemplaza líneas 2844–2869 de Docs/Technical_Report.tex.
% Decisiones que lo respaldan: D-014, D-018, D-026, D-030, D-036, D-037,
% H-011, H-035, H-051, H-054.
% ---

\subsection{Módulo 7: Clasificación}
\label{sec:modulo7}

El módulo de clasificación entrena el mismo perceptrón multicapa en PyTorch sobre la
representación de cada condición: una capa de entrada de dimensión 12, capas ocultas de 32 y 16
neuronas con activación ReLU, y una salida sigmoide que da la probabilidad de malignidad. La pérdida
es la entropía cruzada binaria ponderada (ecuación~\eqref{eq:weighted_bce}), con los pesos de la
ecuación~\eqref{eq:class_weights}, optimizada con Adam (tasa $10^{-3}$, decaimiento de pesos
$10^{-4}$, lotes de 128). Las entradas se estandarizan con la media y la desviación del
entrenamiento de cada ajuste: las representaciones difieren en escala en más de un orden de magnitud (en masas, desviación típica media por dimensión de 0.90 en C1 frente a 0.056 en C4), y
sin estandarizar el optimizador favorecería a las de escala grande por una razón ajena a su
información. El entrenamiento se detiene cuando la pérdida, medida sobre el 15\,\% de los pacientes
del entrenamiento, deja de mejorar durante 40 épocas, y se conservan los mejores pesos. Los
hiperparámetros se fijaron antes de ver los resultados y son idénticos en las cinco condiciones, con
la misma semilla.

Cada condición se evalúa con validación cruzada de cinco \textit{folds}, los fijados en M4, y con un
único ajuste sobre todo el entrenamiento que se evalúa una sola vez en el conjunto de prueba oficial.
La incertidumbre del conjunto de prueba se estima con 2\,000 remuestreos \textit{bootstrap} por
paciente, pareados entre condiciones. El selector y el escalador de M4 se ajustan una sola vez con todo
el entrenamiento, no dentro de cada \textit{fold}, así que la validación cruzada es algo optimista;
volver a seleccionar dentro de cada \textit{fold} cambia el AUC en 0.006 o menos.

Los parámetros $\boldsymbol\theta$ del circuito no se entrenan. La razón principal es de validez
interna: C2 y C3 son transformaciones que no ven las etiquetas, y si C4 y C5 ajustaran
$\boldsymbol\theta$ contra ellas, cualquier ventaja podría deberse al ajuste supervisado y no a la
codificación. Además, la pregunta de investigación se refiere a la codificación $\phi$ y no a un
procedimiento de optimización. Una tercera razón es el costo: con el método de
\textit{parameter-shift}, el que puede usarse en un computador cuántico, entrenar $\boldsymbol\theta$, medido en el simulador,
costaría unas 205 horas por \textit{feature map}. En el simulador existe un atajo: retropropagar a
través del vector de estado da el mismo gradiente unas 3\,600 veces más rápido. Pero ese atajo no
existe en el hardware, porque exige conocer todas las amplitudes del estado. La sensibilidad de la
conclusión a la semilla, la red, el escalado, la parada temprana, \texttt{reps} y $k$ se reporta en
la sección de resultados.
